%!TEX root = ../main.tex

\section{Quesito 7}

\subsection*{Una domanda con importanti implicazioni pratiche (opzionale)}

\vspace{0.5em}

L'approccio \emph{wall function} impiegato in questo caso di studio richiede che le celle vicine alla parete siano caratterizzate da un $y^+$ superiore a 30 e inferiore a 130. Cosa succede se questa condizione non viene soddisfatta? Si noti che, in pratica, in una geometria complessa questa condizione non sarà soddisfatta per ogni cella.

Eseguire le simulazioni numeriche con il modello di turbolenza $k$-$\varepsilon$ standard ignorando intenzionalmente questo vincolo, ovvero calcolare la soluzione CFD su mesh fini vicino alla parete nel range $y^+ = 1 \div 130$. Analizzare cosa accade quando $y^+$ è al di fuori del range ammissibile (cioè per $y^+$ inferiore a 30), facendo riferimento a tutte le variabili di output.

\NotionDivider

Applicare la legge logaritmica dove vige quella lineare genera turbolenza dove il flusso dovrebbe essere quasi laminare (sottostrato viscoso), con importanti conseguenze come:
\begin{itemize}
    \item senza smorzamento (dumping functions), il codice calcolerà una produzione di turbolenza enorme e irreale attaccata al muro;
    \item friction factor sbagliato;
    \item profili di velocità distorti vicino alla parete;
    \item Convergenza fifficile perché le equazioni sono instabili.
\end{itemize}